\begin{table}[htbp]
\centering
\caption{OD-level response to P3 in the Jiangxi corridor}\label{tab:jiangxi_od_outcomes}
\begingroup
\footnotesize
\setlength{\tabcolsep}{0.50em}
\renewcommand{\arraystretch}{1.08}
\begin{tabular}{@{}lrrrr@{}}
\toprule
OD & Profit rail trips (000/day) & Profit rail share (\%) & Welfare rail trips (000/day) & Welfare rail share (\%) \\
\midrule
GZ--JA & 22.6 $\rightarrow$ 21.1 & 80.6 $\rightarrow$ 75.0 & 22.6 $\rightarrow$ 23.0 & 80.6 $\rightarrow$ 81.9 \\
\textbf{GZ--NC} & 9.7 $\rightarrow$ 9.1 & 89.7 $\rightarrow$ 87.9 & 9.7 $\rightarrow$ 12.1 & 89.7 $\rightarrow$ 93.5 \\
GZ--JJ & 0.7 $\rightarrow$ 0.9 & 98.6 $\rightarrow$ 98.9 & 0.7 $\rightarrow$ 1.3 & 98.6 $\rightarrow$ 99.4 \\
JA--NC & 18.6 $\rightarrow$ 17.8 & 81.0 $\rightarrow$ 77.6 & 18.6 $\rightarrow$ 18.9 & 81.0 $\rightarrow$ 82.4 \\
JA--JJ & 5.1 $\rightarrow$ 5.1 & 97.9 $\rightarrow$ 97.5 & 5.1 $\rightarrow$ 5.2 & 97.9 $\rightarrow$ 98.3 \\
NC--JJ & 17.6 $\rightarrow$ 15.1 & 35.4 $\rightarrow$ 30.3 & 17.6 $\rightarrow$ 17.3 & 35.4 $\rightarrow$ 34.7 \\
\bottomrule
\end{tabular}
\parbox{0.96\textwidth}{\footnotesize\textit{Notes.} Each entry reports P0 $\rightarrow$ P3. GZ, JA, NC, and JJ denote Ganzhou, Ji'an, Nanchang, and Jiujiang. The bold GZ--NC row is the required Ganzhou--Nanchang OD.}
\endgroup
\end{table}
